\documentclass[12pt,spanish,mexico]{article}
\usepackage[utf8]{inputenc}
\PassOptionsToPackage{hyphens}{url}
\usepackage{hyperref}
\usepackage[es-nodecimaldot]{babel} % nombre de opción corregido
\usepackage{amsthm, amssymb, amsmath, amsfonts, bbm, mathtools}
\usepackage{csquotes}
\usepackage{enumerate}
\usepackage[shortlabels]{enumitem}
\usepackage{titling}
\usepackage{blindtext}
\usepackage{lmodern} 
\usepackage{float}
\usepackage{mathrsfs}
\usepackage[capitalise]{cleveref}
\usepackage{caption} 
\usepackage{graphicx}
\usepackage{chngcntr}
\usepackage{enumitem}    
\usepackage{cleveref} 
\usepackage[ruled,vlined]{algorithm2e}

\counterwithin{figure}{section}
\numberwithin{equation}{section}
\usepackage[vmargin=1in, hmargin=0.5in]{geometry}  

\setlength{\topmargin}{-50pt}
\setlength{\headheight}{38.23112pt}
\setlength{\headsep}{24pt}

\makeatletter
\def\iftitle{\ifx\thetitle \empty \else }
\def\safetitle{\iftitle{\thetitle}\fi}
\def\headertitle{%
    \ifdefined\hw{\hw}\fi%
    \ifdefined\hw \iftitle{ -- }\fi%
    \safetitle%
}
\def\@maketitle{%
  \begin{center}%
Proyecto: Asimilación de datos 4D-Var con el modelo logístico \\
Muñoz Macías Christian Jaime 
    \vskip 1em%
    \hrule
  \end{center}%
  \par
  \vskip 1.5em%
}
\makeatother
\usepackage{amsmath}
\usepackage{tikz}

\newcommand\RADot{%
  \mathord{%
    \mspace{1mu}%
    \text{\radot}%
    \mspace{1mu}%
  }%
}

\newcommand\radot{%
    \tikz[line cap=round,x=1ex,y=1ex,line width=0.3pt]
    {\draw (0,1) |- (1,0) (0.55,0) arc(0:90:0.55); \fill (0.23,0.23) circle (0.05);}%
}

%% encabezado personalizado
\usepackage{fancyhdr}
%% estilo plain personalizado
\fancypagestyle{plain}{%
    \fancyhf{} % limpia encabezado y pie de página
    \renewcommand{\headrulewidth}{0pt}
    \setlength{\headheight}{0pt}
    \fancyfoot[R]{\thepage}
}
%% estilo para todas las páginas
\pagestyle{fancy}
\fancyhf{}
\fancyhead[C]{Maestría en Matemáticas Aplicadas}
\fancyfoot[R]{\thepage}

%%%%%%%%%%%%%%%%%%%%%%%%
%% DEFINICIONES ÚTILES %%
%%%%%%%%%%%%%%%%%%%%%%%%
%% definiciones matemáticas
\newcommand{\N}{\mathbb{N}}
\newcommand{\Z}{\mathbb{Z}}
\newcommand{\Ha}{\mathcal{H}}
\newcommand{\R}{\mathbb{R}}
\newcommand{\Q}{\mathbb{Q}}
\newcommand{\C}{\mathbb{C}}
\newcommand{\PP}{\mathbb{P}}

\newtheorem{proposition}{Proposición}
\newtheorem{theorem}{Teorema}
\newtheorem{lemma}{Lema}
\newtheorem{corollary}{Corolario}

\theoremstyle{definition}
\newtheorem{problem}{Problema}
\newtheorem{definition}{Definición}[section]
\newtheorem*{example}{Ejemplo}

\theoremstyle{remark}
\newtheorem*{observation}{Observación}
\newtheorem*{note}{Nota}

\newenvironment{solution}{\begin{proof}[Solución]}{\end{proof}}
\newenvironment{dem}{\begin{proof}[Demostración]}{\end{proof}}
\renewcommand{\qedsymbol}{$\blacksquare$}

%% comandos útiles
\newcommand{\eps}{\varepsilon}
\newcommand{\To}{\longrightarrow}
\newcommand{\Eu}{\varPhi}
\newcommand{\Mapsto}{\longmapsto}

%% delimitadores emparejados
\DeclarePairedDelimiter{\set}{\{}{\}}
\DeclarePairedDelimiter{\floor}\lfloor\rfloor
\DeclarePairedDelimiter{\ceil}\lceil\rceil
\DeclarePairedDelimiter{\abs}\lvert\rvert
\DeclarePairedDelimiter{\gen}\langle\rangle
\DeclareMathOperator{\sgn}{sgn}
\DeclareMathOperator{\ord}{ord}

%% operador
\newcommand{\op}[1]{\operatorname{#1}}

%% etiquetas del documento
\title{Proyecto: Asimilación de datos 4D-Var con el modelo logístico}
\def\course{Materia}
\def\university{CIMAT - MsC. Maestría en matemáticas aplicadas}
\author{Muñoz Macías Christian Jaime}
\def\email{\url{christian.munoz@cimat.mx}}
\date{\today}

%% permitir saltos de línea en ecuaciones
\allowdisplaybreaks[4]

\begin{document}

\maketitle

\section*{Introducción}
La asimilación de datos consiste en reconstruir el estado inicial de un sistema dinámico a partir de una secuencia de observaciones ruidosas distribuidas a lo largo de una ventana temporal. Esta es la tarea que resuelven los métodos de asimilación variacional —3D-Var y 4D-Var— desarrollados originalmente en meteorología y oceanografía, y hoy aplicados también en epidemiología, glaciología, sismología y modelado de calidad del aire. En este proyecto implementaremos 4D-Var sobre un modelo sencillo pero suficientemente rico para exhibir todos los ingredientes del método: ecuación forward no lineal, ecuación adjunta lineal retrógrada, gradiente reducido, y término de background con covarianza.

\section{Modelo dinámico}
La dinámica del estado \(x(t)\) está gobernada por la ecuación logística
\[
\frac{dx}{dt} = r x \left(1 - \frac{x}{K}\right), \quad t \in [0, T], \qquad (1)
\]
con condición inicial
\[
x(0) = X_0, \qquad (2)
\]
donde los parámetros \(r > 0\) y \(K > 0\) se suponen conocidos y la incógnita es el estado inicial \(X_0\). Para hacer explícita la dependencia del control, denotamos por \(x(t; X_0)\) la solución de (1)–(2), y abreviamos \(x^\dagger(t) := x(t; X_0^\dagger)\) para la trayectoria verdadera. Disponemos de \(N\) observaciones ruidosas
\[
y_i = x^\dagger(t_i) + \eta_i, \quad \eta_i \sim \mathcal{N}(0, \sigma^2), \quad i = 1, \dots, N, \qquad (3)
\]
en tiempos \(0 < t_1 < \dots < t_N \leq T\) distribuidos a lo largo de la ventana de asimilación \([0, T]\). Adicionalmente contamos con un valor de background \(X_0^b\) —una estimación previa de la condición inicial— con incertidumbre caracterizada por la varianza \(B > 0\).

\section{Funcional 4D-Var}
El método 4D-Var formula la asimilación como el problema de optimización
\[
\min_{X_0 \in \mathbb{R}} J(X_0) = \underbrace{\frac{1}{2B}(X_0 - X_0^b)^2}_{\text{término de background } J_b} + \underbrace{\frac{1}{2\sigma^2} \sum_{i=1}^N (x(t_i; X_0) - y_i)^2}_{\text{término de observaciones } J_o} \qquad (4)
\]
sujeto a la restricción dinámica (1)–(2). La solución \(\hat{X}_0\) de (4) es el análisis; la trayectoria \(x(t; \hat{X}_0)\) obtenida al integrar el modelo desde \(\hat{X}_0\) es el análisis 4D. Bayesianamente, \(\hat{X}_0\) es el estimador de máximo a posteriori bajo verosimilitud gaussiana de los errores y prior gaussiana \(X_0 \sim \mathcal{N}(X_0^b, B)\).

\subsection{Sobre la dimensión}
Aunque en este proyecto el estado es escalar (\(x(t) \in \mathbb{R}\)) y por tanto la covarianza \(B\) es también escalar, todo el formalismo aquí desarrollado es idéntico al caso vectorial, donde \(x(t) \in \mathbb{R}^n\), \(B \in \mathbb{R}^{n \times n}\) es una matriz simétrica positiva definida, y la ecuación adjunta es un sistema de EDOs lineales retrógradas. Toda la maquinaria del adjunto escala al caso vectorial sin cambios estructurales.

\section{Marco teórico}
El problema (4) es una instancia del esquema general del Corolario 3.56 del tutorial de Bui-Thanh (2023), con variable de control \(X_0\) y restricción dada por la EDO. El gradiente \(dJ/dX_0\) se calcula vía el Lagrangiano
\[
\mathcal{L}(x, X_0, \lambda) = J(X_0) + \int_0^T \lambda(t) \left[ \dot{x} - r x(1 - x/K) \right] dt + \mu(x(0) - X_0), \quad (5)
\]
donde \(\lambda(t)\) es el multiplicador para la EDO y \(\mu \in \mathbb{R}\) para la condición inicial.

\section{Datos del problema}
Los parámetros son como se indica en la siguiente tabla:

\begin{table}[H]
\centering
\begin{tabular}{|l|l|c|}
\hline
Nombre & Significado & Valor \\
\hline
\(r\) & Tasa de crecimiento (conocida) & 1.0 \\
\(K\) & Capacidad de carga (conocida) & 100 \\
\(T\) & Horizonte / ventana de asimilación & 5 \\
\(X_0^{\dagger}\) & Estado inicial verdadero (oculto) & 10 \\
\(X_0^{b}\) & Estado inicial background & por elegir \\
\(B\) & Varianza del background & por elegir \\
\(N\) & Número de observaciones en \([0, T]\) & 10 \\
\(\sigma\) & Desviación estándar del ruido & \(0.05 \cdot \max_{t\in[0,T]} x^{\dagger}(t)\) \\
\(\tau_k\) & Paso de descenso en iteración \(k\) & búsqueda de línea \\
\hline
\end{tabular}
\end{table}

\textbf{Observación importante sobre la ventana.} Se sugiere \(T = 5\). Con \(r = 1\) y \(X_0^{\dagger} = 10\), el sistema satura en \(t \approx 7\); tomar la ventana hasta \(t = 5\) asegura que las observaciones contengan información sobre la fase de crecimiento, donde \(X_0\) es realmente identificable. Discute en tu reporte qué ocurre si se toma \(T\) demasiado grande (saturación) o demasiado pequeño (poca información).

\section{Resolución del problema de asimilación 4D-Var}

\subsection{1. Deducción de la ecuación adjunta}
Partiendo del Lagrangiano (5), calcula la variación \(\langle \partial\mathcal{L}/\partial x, h \rangle\) en una dirección arbitraria \(h(t)\) con \(h(0)\) libre, integra por partes y anula la variación. Obtén la ecuación adjunta
\[
-\dot{\lambda} - r \left( 1 - \frac{2x(t; X_0)}{K} \right) \lambda = -\frac{1}{\sigma^2} \sum_{i=1}^{N} \left( x(t_i; X_0) - y_i \right) \delta(t - t_i), \quad \lambda(T) = 0, \quad (6)
\]
junto con la condición de transversalidad \(\mu = \lambda(0)\) que relaciona el multiplicador de la EDO con el de la condición inicial. Aquí \(x(t; X_0)\) es la solución forward para el valor actual del control, es decir, el coeficiente de la EDO adjunta depende de la trayectoria forward y debe almacenarse durante la simulación del problema directo.

\begin{solution}
El funcional es:
\[
J(X_0) = \frac{1}{2B}(X_0 - X_0^b)^2 + \frac{1}{2\sigma^2} \sum_{i=1}^N (x(t_i; X_0) - y_i)^2
\]
Tenemos como Lagrangiano:
\[
\mathcal{L}(x, X_0, \lambda) = J(X_0) + \int_0^T \lambda(t) \left[ \dot{x} - r x(1 - x/K) \right] dt + \mu(x(0) - X_0), 
\]
sabemos que minimizar tiene que cumplir que: 
$\nabla J= \nabla \mathcal{L}$
entonces la derivada direccional de $J$ es:
\begin{align*}
D J(X_0)h =& D \mathcal{L}(x,X_0)(\xi,h) \\
=& D_x \mathcal{L} \xi+D_{X_0} \mathcal{L} h
\end{align*}


Lo anterior es válido porque en el punto estacionario las variaciones respecto a $\lambda$ se anulan (recuperan la ecuación forward), de modo que $D\mathcal{L}$ colapsa a $DJ$.

derivando respecto a $x$:
\[
D_x \mathcal{L} \xi = \frac{1}{\sigma^2} \sum_{i=1}^N (x(t_i) - y_i) \xi(t_i)+  \int_0^T \lambda(t) \left[ \dot{\xi}-r\left(1 -\frac{2x}{K} \right)\xi \right] dt + \mu\,\xi(0),
\]

intercambiar el diferencial con la integral es válido por la regularidad de $\lambda$ y $\xi$.

derivando respecto a $X_0$:
\[
D_{X_0} \mathcal{L} h= \frac{1}{B}(X_0 - X_0^b)h - \mu h, 
\]

nos interesa que $D_x \mathcal{L} \xi=0$, para ello trabajamos el término:
\[
\int_0^T \lambda(t) \left[ \dot{\xi}-r\left(1 -\frac{2x}{K} \right)\xi \right] dt 
\]
usando la regla del producto en $\lambda(t)\dot{\xi}$:
\[
\frac{d}{dt}(\lambda \xi)=\dot{\lambda} \xi+\lambda \dot{\xi}
\implies
\lambda \dot{\xi}=\frac{d}{dt}(\lambda \xi)-\dot{\lambda} \xi
\]
entonces la integral queda:
\[
\int_0^T \frac{d}{dt}(\lambda \xi)-\dot{\lambda} \xi-r\left(1 -\frac{2x}{K} \right)\lambda \xi \;  dt 
\]
evaluando el término de frontera:
\[
\int_0^T \frac{d}{dt}(\lambda \xi)\, dt=\lambda(T)\xi(T)-\lambda(0)\xi(0)=-\lambda(0)\xi(0)
\]
donde usamos $\lambda(T)=0$. Sustituyendo en $D_x\mathcal{L}\,\xi = 0$:
\[
\frac{1}{\sigma^2} \sum_{i=1}^N (x(t_i) - y_i) \xi(t_i)+\int_0^T\left[-\dot{\lambda} -r\left(1 -\frac{2x}{K} \right)\lambda\right] \xi \;  dt + (\mu - \lambda(0))\,\xi(0)= 0, 
\]
como $\xi(0)$ es libre, para que esto se anule para toda dirección $\xi$ el coeficiente de $\xi(0)$ debe ser cero:
\[
\mu = \lambda(0) \qquad \text{(condición de transversalidad)}
\]
cancelado ese término, usando que $\xi(t_i)=\int_0^T \xi(t)\,\delta(t-t_i)\,dt$, la condición debe valer para todo $\xi$, lo que implica la igualdad puntual:
\[
-\dot{\lambda}-r\left(1 -\frac{2x}{K} \right)\lambda=- \frac{1}{\sigma^2} \sum_{i=1}^N (x(t_i) - y_i) \delta (t-t_i)
\]
hemos deducido la ecuación adjunta con condición terminal $\lambda(T)=0$.
\end{solution}
\subsection{2. Deducción del gradiente reducido}
Varía \(\mathcal{L}\) respecto a \(X_0\), usa que \(x, \lambda\) satisfacen las ecuaciones forward y adjunta, y deduce
\[
\frac{dJ}{dX_0} = \frac{X_0 - X_0^b}{B} - \lambda(0). \quad (7)
\]
Comenta el significado: el gradiente es la suma de un tirón hacia el background (primer término) y la sensibilidad retrógrada de las observaciones al estado inicial (segundo término, \(\lambda(0)\)). Esto justifica llamar al adjunto el “transportador hacia atrás de la información de las observaciones”.

\begin{solution}
Variamos $\mathcal{L}$ respecto a $X_0$ en una dirección $h$:
\[
D_{X_0} \mathcal{L}\, h = \frac{1}{B}(X_0 - X_0^b)h - \mu h
\]
como $x$ y $\lambda$ satisfacen las ecuaciones forward y adjunta, esta variación coincide con $DJ(X_0)h$, entonces:
\[
\frac{dJ}{dX_0} h= \frac{1}{B}(X_0 - X_0^b)h - \mu h
\]
sustituyendo la condición de transversalidad $\mu = \lambda(0)$ obtenida en el punto anterior:
\[
\boxed{\frac{dJ}{dX_0} = \frac{X_0 - X_0^b}{B} - \lambda(0)}
\]
El primer término jala la solución hacia el background; el segundo, $-\lambda(0)$, transporta hacia atrás la información de las observaciones a través de la ecuación adjunta.
\\
\\
El gradiente tiene una interpretación clara: si se busca minimizar $J$, el término 
$\frac{X_0 - X_0^b}{B}$ empuja a $X_0$ hacia el background $X_0^b$, pues se anula 
exactamente cuando $X_0 = X_0^b$. Por otro lado, $\lambda(0)$ es la solución en $t=0$ 
de la ecuación adjunta (6), que es una EDO lineal de primer orden no homogénea retrógrada 
en $[0,T]$: su término forzante son las discrepancias $(x(t_i;X_0)-y_i)$ en los tiempos 
de observación, y al integrar de $T$ hacia $0$ propaga esa información hasta el estado 
inicial. De ahí que $-\lambda(0)$ mida la sensibilidad acumulada de todas las observaciones 
respecto a $X_0$, jalando el análisis hacia donde los datos son más consistentes.

\end{solution}

\subsection{3. Generación de datos sintéticos (twin experiment)}
Adopta el protocolo estándar de validación 4D-Var:
\begin{enumerate}[label=(\alph*)]
\item Integra (1) con \(X_0^\dagger = 10\) usando RK4 sobre \([0, 5]\);
\being{solution}

\end{solution}
\item Muestra \(N = 10\) observaciones en tiempos equiespaciados \(t_i = 0.5\, i\) para \(i = 1, \dots, 10\);
\item Añade ruido gaussiano con \(\sigma = 0.05 \cdot \max_{t\in[0,T]} x^\dagger(t)\).
\end{enumerate}
La trayectoria verdadera te servirá únicamente como referencia para medir el error del análisis; el algoritmo nunca la verá.

\subsection{4. Implementación forward}
Implementa un integrador RK4 para (1) que reciba \(X_0\) y devuelva la trayectoria \(x(t; X_0)\) sobre una malla densa que contenga los \(t_i\). Almacena la trayectoria; el adjunto la necesitará como coeficiente.

\subsection{5. Implementación del adjunto}
Implementa la integración retrógrada de (6) de \(T\) a \(0\), partiendo de \(\lambda(T) = 0\). Al cruzar cada \(t_i\) de derecha a izquierda, incorpora el salto
\[
\lambda(t_i^-) = \lambda(t_i^+) + \frac{1}{\sigma^2}(x(t_i; X_0) - y_i).
\]
Verifica con cuidado la consistencia de signos. El signo del salto se obtiene integrando la \(\delta\) de Dirac de la ecuación adjunta (6) sobre un entorno \([t_i-\epsilon, t_i+\epsilon]\) y haciendo \(\epsilon \to 0\); con la convención de signos de (6) y la integración de \(T\) hacia \(0\), el incremento \(\lambda(t_i^-) - \lambda(t_i^+)\) es positivo cuando \(x(t_i; X_0) > y_i\). Un error de signo aquí es la causa más frecuente de fallo del test de Taylor (punto siguiente). Devuelve \(\lambda(0)\) y, con él, ensambla el gradiente (7).

\subsection{6. Test de Taylor (opcional)}
Antes de optimizar, valida tu gradiente. Para una dirección \(\delta X_0 \in \mathbb{R}\) arbitraria y para \(\varepsilon \in \{10^{-1}, 10^{-2}, \dots, 10^{-8}\}\), calcula
\[
E_1(\varepsilon) = |J(X_0 + \varepsilon \delta X_0) - J(X_0)|, \quad E_2(\varepsilon) = |J(X_0 + \varepsilon \delta X_0) - (J(X_0) + \varepsilon \left(\frac{dJ}{dX_0}\right) \delta X_0)|.
\]
\(E_1\) debe decaer con pendiente 1 y \(E_2\) con pendiente 2 en gráfica log-log, ambos sobre al menos cuatro órdenes de magnitud. Si la pendiente de \(E_2\) no es 2, el adjunto está mal y debes corregirlo antes de continuar.

\subsection{7. Minimización 4D-Var}
Implementa descenso más pronunciado con búsqueda de línea de Armijo:
\begin{enumerate}[label=(\alph*)]
\item Resolver forward con \(X_0^{(k)}\) para obtener \(x_k(t)\);
\item Resolver adjunto para obtener \(\lambda_k(0)\);
\item Ensamblar el gradiente (7);
\item Línea de búsqueda y actualizar \(X_0^{(k+1)} = X_0^{(k)} - \tau_k \frac{dJ}{dX_0}\).
\end{enumerate}
Inicializa con \(X_0^{(0)} = X_0^b\) (el background). Reporta la convergencia de \(J(X_0^{(k)})\), \(|dJ/dX_0|\) y \(|X_0^{(k)} - X_0^\dagger|\).

\subsection{8. Estudio del peso relativo background/observaciones}
Realiza tres experimentos de asimilación variando la confianza relativa entre el background y las observaciones, manteniendo \(X_0^b = 5\) (lejos del valor verdadero 10):
\begin{enumerate}[label=(\alph*)]
\item Background dominante: \(B = 0.1\), \(\sigma\) como en la tabla;
\item Equilibrio: \(B = 10\);
\item Observaciones dominantes: \(B = 10^{4}\).
\end{enumerate}
En cada caso reporta el análisis \(\hat{X}_0\) y compáralo con \(X_0^\dagger = 10\) y con \(X_0^b = 5\). Discute cómo \(B\) controla el compromiso entre confiar en el prior y confiar en los datos.

\section{Entregables}
\begin{enumerate}[label=(\alph*)]
\item Reporte en PDF (máximo 12 páginas) con las deducciones de los puntos 1-2, la figura obligatoria del test de Taylor (punto 6), y los experimentos de los puntos 7-11.
\item Código reproducible (Python, Julia o MATLAB) con un README que permita regenerar todas las figuras del reporte.
\end{enumerate}

\section*{Referencias sugeridas}
\begin{itemize}
\item Bui-Thanh, T. (2023). Adjoint and Its roles in Sciences, Engineering, and Mathematics: A Tutorial. arXiv:2306.09917. Secciones 3.5-3.6.
\item Asch, M., Bocquet, M., Nodet, M. (2016). Data Assimilation: Methods, Algorithms, and Applications. SIAM. Capítulos 2 y 4.
\item Lewis, J.M., Lakshmivarahan, S., Dhall, S. (2006). Dynamic Data Assimilation: A Least Squares Approach. Cambridge University Press.
\item Le Dimet, F.-X., Talagrand, O. (1986). Variational algorithms for analysis and assimilation of meteorological observations. Tellus A, 38(2), 97-110. (Artículo fundacional de 4D-Var.)
\end{itemize}

\end{document}